% PROJECT-SPECIFIC SUPPORTING STATEMENT TEMPLATE
% Upload to Overleaf as main.tex. Compile with pdfLaTeX.
% Replace all bracketed text with your own prose.
% Comments beginning with % will not appear in the PDF.
% Follow any word/page limit or prompt given in the application portal.
% The paragraph lengths below are optional drafting guides, not requirements.
% Use connected paragraphs; the prompts are not submission-ready prose.
%
% Project advert checked 15 October 2026:
% https://www.findaphd.com/phds/project/?p198781

\documentclass[11pt,a4paper]{article}
\usepackage[margin=1in]{geometry}
\usepackage[T1]{fontenc}
\usepackage[utf8]{inputenc}
\usepackage{lmodern}
\usepackage{microtype}
\usepackage{setspace}
\usepackage{fancyhdr}
\usepackage[hidelinks]{hyperref}

% ---------- Application details ----------
\newcommand{\ApplicantName}{Bahar Barghbani}
\newcommand{\UniversityName}{The University of Manchester}
\newcommand{\ProjectTitle}{Resource-Efficient Execution of Data Centre-Grade
Workloads with Lightweight Operating Systems on Emerging Capability Hardware}
\newcommand{\Supervisors}{Dr Pierre Olivier and Dr Pavlos Petoumenos}

% ---------- Formatting ----------
\setstretch{1.08}
\setlength{\parindent}{0pt}
\setlength{\parskip}{0.7em}
\setlength{\emergencystretch}{2em}
\pagestyle{fancy}
\fancyhf{}
\renewcommand{\headrulewidth}{0pt}
\fancyfoot[C]{\small\thepage}
\hypersetup{pdftitle={Supporting Statement},pdfauthor={Bahar Barghbani}}

\begin{document}

\begin{center}
    {\Large\bfseries Supporting Statement}\par
    \vspace{0.35em}
    {\normalsize\bfseries\ApplicantName}\par
    \vspace{0.25em}
    {\small\UniversityName}\par
    \vspace{0.5em}
    \begin{minipage}{0.95\textwidth}
        \centering\small\ProjectTitle
    \end{minipage}\par
    \vspace{0.35em}
    {\small Supervisors: \Supervisors}
\end{center}
\vspace{0.3em}

I became interested in computer systems during my third semester at Sharif University of Technology. In a course on computer structure and languages, I learned about instruction set architectures and assembly programming. I enjoyed being able to follow what happened beneath the code I wrote. Later, my Computer Architecture and Operating Systems courses helped me decide that I wanted to continue studying systems, particularly processor architecture and memory behaviour. I am now a final-year Computer Engineering student applying for the PhD project “AI-driven Hardware-Software Co-optimisation for Sustainable and Efficient Computing” at the University of Manchester.

My interest in understanding how things work also comes from studying physics. I earned a bronze medal in Iran’s National Physics Olympiad, and I think part of what attracts me to architecture is the ability to connect a program’s behaviour to the hardware executing it. Studying caches and the memory hierarchy gave me a more specific area to explore, and last winter I began working with Professor Alberto Ros at the University of Murcia on instruction-cache prefetching.

Last winter, I began working with Professor Alberto Ros at the University of Murcia on instruction-cache prefetching. My initial task was to reproduce published results using a gem5 implementation of the Entangling Instruction Prefetcher. When the IPC numbers did not match, I inspected the code and configuration files and discussed possible causes with a researcher who had worked on the implementation. After several revisions, my baseline results matched the reference values for two workloads, although the prefetcher results still differed. The researcher subsequently observed variation in his own runs as well. We later moved to an extended version of ChampSim to continue the research.

My current work asks whether some instruction prefetches should fill L2 instead of L1I. The motivation is to investigate whether selective placement can reduce L1I pollution enough to compensate for the additional delay when an instruction is later needed. Following Professor Ros’s suggestions, I tested different placement choices and implemented a policy using Entangling’s existing confidence counters. In this policy, lower-confidence predictions go to L2, while higher-confidence predictions continue to fill L1I. I also tested a variation that keeps the first line of a predicted block in L1I and applies the confidence-based choice only to its continuation lines.

Changing the destination also required changes to the prefetcher’s tracking. The original implementation expected tracked requests to complete with an L1I fill, so requests that stopped in L2 could leave entries occupied indefinitely. I introduced separate tracking for L2-targeted requests and extended the simulator to retain their identity through request merges and report their fills and evictions. This allowed me to preserve information about when a prefetch arrived and whether the line was subsequently used.

I then compared two ways of using this information: reconstructing timing information for learning source–destination relationships, and updating a relationship’s confidence when an L2-prefetched line was used or evicted unused. I evaluated the mechanisms separately and together on two CVP traces. Confidence feedback alone produced the highest IPC among the tested L2 variants, but neither mechanism brought them above the original all-L1I configuration. These results have made the choice of placement signal a central question in my work: a counter that measures confidence in a prediction may not tell us which cache level should receive it.
I am currently using CVP workload traces to investigate where prefetched instruction blocks should be placed in the cache hierarchy. I am testing policies that introduce L2 prefetching and comparing them with a configuration in which all prefetched instructions fill the L1 instruction cache. My aim is to determine whether a different placement policy can improve performance on these traces. The work is still ongoing. It has given me experience with architectural simulation and a research question that I want to continue exploring: what information should guide a prefetching decision, and how should that decision change for different workloads?

During a six-week internship with Professor Özcan Öztürk at Sabancı University, I implemented compute-intensive and memory-intensive benchmarks from scratch. I also studied methods of analysing performance using hardware performance monitoring units (PMUs). This work introduced me to another way of investigating memory behaviour, through measurements collected while programs run on real hardware. We plan to extend the work by investigating whether metrics derived from PMU data can help characterise performance. That part of the research is still at the preparatory stage.

The Manchester project interests me because it offers a way to bring information from program code into hardware optimisation decisions. In my current research, I work with execution traces and prefetching policies. I would like to understand how information extracted before execution could help those policies make better decisions. The proposed use of LLMs to identify memory-access characteristics in source code and supply hints to cache mechanisms is therefore closely related to the questions I have begun working on. I am particularly interested in whether these hints can reduce the hardware resources needed to learn useful patterns during execution.

My experience so far is mainly in architecture and performance analysis. Through this PhD, I want to learn how to use machine learning for program analysis and how to pass the resulting information to hardware through the compiler and binary. I would also like to extend my experimental work beyond performance comparisons to consider the energy use and hardware cost of an optimisation.

I want to continue researching memory systems, with a particular interest in caches and the processor front end. Working with Dr Pavlos Petoumenos and Dr Pierre Olivier on this project would allow me to build on my prefetching research while learning how software information can inform hardware design. This is the direction I would like to pursue during my doctoral studies.
\end{document}
